% ===== §1 Prelude =====
\section{Prelude: A Ring and a Flower}
\label{sec:prelude}

\noindent
Two scenes.

In the first, a person walks into a jeweller's and purchases a diamond ring. The ring is expensive---the price of several months' salary, perhaps more. It is beautifully made: a round brilliant cut stone in a platinum setting, designed by craftspeople whose skills have been refined across generations, sourced from materials that took the earth millions of years to produce. The purchase is made deliberately, with care. It is intended as an expression of love.

In the second scene, a different person is walking a familiar path---a path walked many times, in different seasons and different moods. On the verge of the path, in a crack between the paving stones, there is a flower. It is not a rare flower. It has no particular monetary value. But the person stops, and in the moment of stopping, something happens: they think of the person they love---not in general, not as an abstraction, but specifically, in all the particular detail that months or years of shared life have made available. They pick the flower. They bring it home.

\emc{Which of these is the more luxurious gift?}

The answer that the dominant framework of luxury theory would give is obvious: the diamond ring. It is more expensive, more durable, more socially recognised as an expression of serious commitment, more legible within the symbolic order as evidence of the depth of the giver's investment. The flower is charming but ephemeral, cheap, easily obtained. It is, in the terms of the dominant framework, not a luxury at all.

This paper argues that the dominant framework is wrong---not partially wrong, not wrong about some cases, but structurally wrong, missing the dimension of value that is most fundamental to intimate relational life. It argues that the flower, in the right circumstances, is more luxurious than the ring---not because it is rarer or more expensive or more socially prestigious, but because it carries a higher degree of coupling between the object and the relational subject's spatiotemporal structure. And it argues that this coupling degree---not economic value, not sign exchange value, not social prestige---is the true measure of luxury in intimate relations.

\begin{claim}[The central claim]
The true measure of luxury in intimate relations is the degree of coupling between an object and a relational subject's spatiotemporal structure: the extent to which a subject's existence---their intentionality, their presence, their choice made specifically for this relational other at this specific moment---is embedded within the object. Economic value and coupling degree are orthogonal dimensions: a flower can be more luxurious than a diamond, and a diamond can be more luxurious than a flower, depending on the degree of subject-embedding each carries.
\end{claim}

This claim has consequences that extend beyond the question of what makes a good gift. It reveals a systematic misalignment between the dominant framework of luxury---organised around economic value, sign exchange value, and social prestige---and the relational reality of intimate life. It identifies a form of injustice specific to the domain of intimate relational luxury: the misrecognition of subject-embedding, the substitution of economic value for existential presence, the collective confusion of what luxury is. And it generates a justice theory of choice: the claim that the selection of relational luxury over economic luxury is a justice stance---a refusal to use exchangeable symbols to bear what is inchangeable, a commitment to the primacy of the real register over the symbolic.

\emc{A word on this paper's place in the series.} \pinkref{Paper~XIV} established that happiness in intimate relations is the emergent signal of co-evolutionary activity in the relational field \citep{paperxiv}. \pinkref{Paper~XV} developed the political economy of that field, showing that Generative Relational Wealth (GRW) is its primary form of wealth---inexhaustible under use, non-possessable by either party, stored in the geometry of the shared attractor landscape \citep{paperxv}. The present paper asks: what objects can serve as the material media of this wealth? What physical things can carry, in their materiality, the trace of co-evolutionary events, the mark of shared existence, the existential embedding of one subject in the life of another? The answer is: relational luxury objects---objects whose coupling degree with a relational subject's spatiotemporal structure is high. The diamond ring and the flower are the paper's paradigm cases, but the analysis extends to every object that might, in the right circumstances, carry the mark of genuine subject-embedding.

\emc{The paper proceeds through three doctrines and a justice theory.} The first doctrine (\xref{sec:nature_art}) analyses luxury as the transformation of natural materials through human craft---the Bulgari paradigm, in which the beauty of the natural world is given permanent form through the excellence of human making. The second doctrine (\xref{sec:baudrillard}) follows Baudrillard's critique of the consumer society, showing how luxury functions in the symbolic order as a system of social differentiation \citep{baudrillard1981,baudrillard1998}. The third doctrine (\xref{sec:relational}) proposes the concept of relational luxury, defined by coupling degree, and develops its implications for the understanding of intimate relational life. The synthesis (\xref{sec:synthesis}) shows how the three doctrines interact. The justice theory (\xref{sec:justice}) develops the paper's normative conclusions.
